\subsection{Foliations}

In this section, X denotes a manifold of class C \(^{r}\) , with \(r \in \mathrm{N}_{\mathrm{K}}\) . Unless otherwise stated, all manifolds and morphisms considered are assumed to be of class C \(^{r}\) .

9.2.1. Let S be a manifold, and let \(p: X \to S\) be a submersion. For every \(s \in S\) , we endow \(X_s = p^{-1}(s)\) with the manifold structure induced by that of X (5.10.5); the set X is the disjoint union of the \(X_s\) . Denote by \(X_p\) the manifold structure on X obtained by gluing the \(X_s\) , for \(s \in S\) (5.2.4); it is the unique manifold structure on X for which the \(X_s\) are open submanifolds. The topological space \(X_p\) is the sum of the spaces \(X_s\) .

Let V be a manifold and \(f\) a map from V into X. For \(f\) to be a morphism from V to \(\mathbf{X}_{p}\), it is necessary and sufficient that \(f\) be a morphism from V to X and that \(p \circ f\) be locally constant.

9.2.2. A foliation of X is a manifold Y having the same underlying set as X and satisfying the following condition:

\begin{enumerate}
\def\labelenumi{(\Alph{enumi})}
\setcounter{enumi}{5}
\tightlist
\item
  For every \(x \in X\) , there exists an open submanifold U of X containing x, a manifold S, and a submersion \(p: U \to S\) such that the manifold \(U_{p}\) is an open submanifold of Y.
\end{enumerate}

We also say that the pair (X, Y) is a foliated manifold. If (X, Y) and (X', Y') are foliated manifolds, a morphism from (X, Y) to (X', Y') is any map \(f: X \to X'\) which is both a morphism from the manifold X to the manifold X' and a morphism from the manifold Y to the manifold Y'.

Let (X, Y) be a foliated manifold. The identity map Y → X is a bijective immersion. A subset U of X is called a leaf if it is an open set of Y; in this case, U is endowed with the topological structure and the manifold structure induced by those of Y; for example, U is said to be a connected leaf if it is an open and connected subset of Y. The leaves that are submanifolds of X form a basis for the topology of Y. When K = R or C, the connected components of Y are leaves, called the maximal connected leaves of (X, Y).

9.2.3. If \(s \in \mathbf{N}_{\mathbb{K}}\) , \(s \leqslant r\) , we call a foliation of class \(\mathbf{C}^s\) of \(\mathbf{X}\) a foliation of the manifold of class \(\mathbf{C}^s\) underlying \(\mathbf{X}\) .

9.2.4. Examples

\begin{enumerate}
\def\labelenumi{\alph{enumi})}
\item
  The manifold X is a foliation of X, called the coarse foliation of X.
\item
  If \(p: \mathbf{X} \to \mathbf{S}\) is a submersion, \(\mathbf{X}_p\) is a foliation of \(\mathbf{X}\) called the foliation of \(\mathbf{X}\) defined by \(p\) . The foliation of \(\mathbf{X}\) defined by the identity map is the set \(\mathbf{X}\) endowed with its structure of a pure manifold of dimension 0 (5.2.1); it is called the discrete foliation of \(\mathbf{X}\) .
\item
  Let E be a Banach space, F a closed vector subspace of E admitting a topological complement, and p the canonical projection \(E \rightarrow E/F\) . The foliation \(E_{p}\) of E is called the foliation of E defined by F.
\item
  Let \(\Gamma\) be a discrete group acting properly and freely on the topological space X (TG, III, § 4, no. 4). Let Y be a foliation of X, and suppose that, for every \(s \in \Gamma\) , the map \(x \mapsto sx\) is an automorphism of (X, Y). The equivalence relation on X (resp. Y) whose classes are the orbits of \(\Gamma\) is regular (5.9.5). If we denote by X/\(\Gamma\) (resp. Y/\(\Gamma\)) the corresponding quotient manifold, the pair (X/\(\Gamma\),Y/\(\Gamma\)) is a foliated manifold, called the quotient of (X, Y) by \(\Gamma\) .
\end{enumerate}

9.2.5. Let Y be a foliation of X and U an open subset of X. Then U is open in Y, and U, endowed with the manifold structure induced by that of Y, is a foliation of X, called the foliation induced by Y.

More generally, let \(f: X' \to X\) be a morphism such that \(f\) and the identity map \(Y \to X\) form a transversal pair (5.11.1). The fibered product \(Y' = X' \times_X Y\) is canonically identified with a foliation of \(X'\) called the pullback foliation of \(Y\) by \(f\) .

9.2.6. Let \(p: \mathbf{X} \to \mathbf{S}\) and \(p': \mathbf{X}' \to \mathbf{S}'\) be two submersions, and let \(f\) be a morphism from \(\mathbf{X}\) to \(\mathbf{X}'\) . For \(f\) to be a morphism from \(\mathbf{X}_p\) to \(\mathbf{X}_{p'}\) , it is necessary and sufficient that, for every \(x \in \mathbf{X}\) , there exists an open neighborhood \(\mathbf{U}\) of \(x\) in \(\mathbf{X}\) and a morphism \(g\) from \(p(\mathbf{U})\) to \(\mathbf{S}'\) such that \(g \circ p = p' \circ f\) on \(\mathbf{U}\) .

9.2.7. Let Y be a foliation of X. A foliated chart of (X, Y) is a quadruple (U, \(\varphi\) , E, F) such that (U, \(\varphi\) , E) is a chart of X, F is a closed vector subspace of the Banach space E admitting a topological complement, and \(\varphi\) is an isomorphism from the foliation of U induced by Y to the foliation of \(\varphi(\mathbf{U})\) defined by the canonical mapping from E onto E/F. For every point \(x \in \mathbf{X}\) there exists a foliated chart (U, \(\varphi\), E, F) of (X, Y) such that \(x \in \mathbf{U}\) . Let \(n = \dim_x \mathbf{X}\) and let \(m = \dim_x \mathbf{Y}\) . If \(n\) is finite, there exists an open neighborhood U of \(x\) and a coordinate system \(\zeta^1, \ldots, \zeta^n\) of X over U such that the foliation of U induced by Y coincides with the foliation defined by the map (\(\zeta^{m+1}, \ldots, \zeta^n\)) from U into \(\mathbf{K}^{n-m}\) .

9.2.8. Let Y be a foliation of X. When x ranges over X, the spaces \(T_{x}(Y)\) are the fibers of a vector subbundle of class \(C^{r-1}\) of T(X), which is called the subbundle of T(X) tangent to the foliation Y and is denoted T(X, Y). If the foliation Y is defined by a submersion \(p: X \to S\) , then T(X, Y) = Ker T(p) is the relative tangent bundle T(X/S) of X over S (8.1.3).

Let (X, Y) and (X', Y') be two foliated manifolds, and let f be a morphism from X to X'. For f to be a morphism of foliated manifolds, it is necessary and sufficient that T(f) map T(X, Y) into T(X', Y'). In particular, let \(f: Z \to X\) be a morphism of manifolds; for f to be a morphism from Z to Y, it is necessary and sufficient that T(f) map T(Z) into T(X, Y). For a submanifold Z of X to be a leaf of (X, Y), it is necessary and sufficient that one have

\[
\mathrm{T} _ {z} (\mathrm{Z}) = \mathrm{T} _ {z} (\mathrm{X}, \mathrm{Y}) \quad \text {for all} \quad z \in \mathrm{Z}.
\]

Let Y be a foliation of X and let U be a leaf of Y, admitting a countable base of open sets. Let Z be a manifold and f a map from Z to U. For f to be a morphism of manifolds from Z to U, it is necessary and sufficient that the composite of f and the canonical injection of U into X be a morphism of manifolds from Z to X. If X is locally compact and countable at infinity, every connected leaf of (X, Y) admits a countable base of open sets (cf.~TG, I, 3rd ed., §11, no. 7, cor. 2 of th. 1).

9.2.9. Suppose that K = R or C. Let Y be a foliation of X. The following conditions are equivalent:

\begin{enumerate}
\def\labelenumi{\alph{enumi})}
\item
  There exist a manifold S and a submersion \(p: X \to S\) such that \(Y = X_p\) .
\item
  For every \(x \in X\) , there exists a submanifold \(S_{x}\) of X possessing the following two properties:
\begin{enumerate}
\item[\(b_1)\)] The tangent space \(\mathbf{T}_{x}(\mathbf{S}_{x})\) to \(\mathbf{S}_{x}\) at \(x\) is a topological supplement of \(\mathbf{T}_{x}(\mathbf{X},\mathbf{Y})\) in \(\mathbf{T}_{x}(\mathbf{X})\) .
\item[\(b_2)\)] Every connected leaf of (X, Y) intersects \(S_{x}\) at at most one point.
\end{enumerate}
\end{enumerate}

Suppose these conditions are satisfied, and let R be the relation on X “x and y belong to the same connected leaf”; then R is a regular equivalence relation (5.9.5) on X, and, if p denotes the canonical projection \(X \to X/R\), we have \(Y = X_p\) .

Example. --- Take K = R and X = R². Let the discrete group Γ = Z² act on X by translations. Let m ∈ R and let (X, Yₘ) be the foliation defined by the submersion p: (x₁, x₂) ↦ x₂ - mx₁ from X onto R. Set X' = X/Γ and Yₘ' = Yₘ/Γ; then Yₘ' is a foliation of the torus X' (cf.~9.2.4, example d)). If m is rational, there exists a submersion p' from X' onto R/Z such that Yₘ' = Xₚ'. If m is irrational, every maximal connected leaf of (X', Yₘ') is dense in X', and there does not exist a submersion p' from X' into a manifold S such that Yₘ' = Xₚ'.

9.2.10. Let Y be a foliation of X, and let \(\pi: X \to S\) be a morphism of manifolds such that the composite \(Y \to X \xrightarrow{\pi} S\) is étale. If \(S'\) is an open subset of S, a section \(\sigma: S' \to X\) of \(\pi\) above \(S'\) is said to be horizontal (with respect to the foliation Y) if it is a morphism from \(S'\) into Y; it is equivalent to say that \(\sigma(S')\) is a leaf of (X, Y), or again that \(T(\sigma)\) maps \(T(S')\) to \(T(X, Y)\) . For every \(s_0 \in S\) and for every \(x_0 \in \pi^{-1}(s_0)\) , there exists a horizontal section defined in a neighborhood of \(s_0\) and taking the value \(x_0\) at \(s_0\) ; two such sections coincide on a neighborhood of \(s_0\) .

More generally, let \(s_{0}\in S\) , let T be a manifold, let \(f\) be a morphism from T into the submanifold \(\pi^{-1}(s_{0})\) of X, and let \(t_{0}\in T\) . There then exists an open neighborhood \(\mathbf{T}^{\prime}\) (resp. \(S'\)) of \(t_{0}\) in T (resp. of \(s_{0}\) in S), and a morphism

\[
\mathbf {F}: \mathbf {S} ^ {\prime} \times \mathbf {T} ^ {\prime} \rightarrow \mathbf {X}
\]

satisfying the following condition: for every \(t \in \mathbf{T}'\) , the map \(s \mapsto \mathrm{F}(s, t)\) is a horizontal section of \(\pi\) above \(\mathbf{S}'\) taking the value \(f(t)\) at the point \(s_0\) . Two such maps F coincide on a neighborhood of \((s_0, t_0)\) .

